\documentclass[linenumber]{jdsart}
\usepackage{amsmath,amssymb,booktabs,graphicx,pgfplots}
\pgfplotsset{compat=1.18}
\articletype{research-article}
\nocopyrightline
% The production template also serves submissions. Omit publisher-assigned metadata.
\fancypagestyle{copyright}{\fancyhf{}\fancyhead[R]{\thepage}\renewcommand{\headrulewidth}{0pt}}
\makeatletter
\newcommand{\setval}[2]{\expandafter\def\csname pval@#1\endcsname{#2}}
\newcommand{\val}[1]{\ifcsname pval@#1\endcsname\csname pval@#1\endcsname\else\PackageError{paper}{Missing value #1}{}\fi}
\makeatother
\setval{n-bank_leverage}{20,719}
\setval{n-bank_cash_liquidity}{21,473}
\setval{n-bhc_leverage}{11,323}
\setval{n-bhc_cash_liquidity}{11,353}
\setval{n-eia_plants}{11,583}
\setval{py-version}{3.12.10}
\setval{pandas-version}{3.1.1}
\setval{runtime-bank-min}{27}
\setval{runtime-eia-min}{11}
\setval{b1-bank-range}{0.8--2.6\%}
\setval{b1-bank-tv-range}{0.15--0.41}
\setval{b1-eia}{17.3\%}
\setval{b1-eia-harm}{1.70}
\setval{b1-eia-tv}{0.18}
\setval{surv-bank-range}{15.9--23.8\%}
\setval{surv-eia-range}{42.0--47.7\%}
\setval{b1-harm-sub}{0.57}
\setval{b1-harm-full}{0.90}
\setval{b1-rel-full-bl}{1.7\%}
\setval{b1-bank-tv-max}{41\%}
\setval{t3-max-sel}{0}
\setval{rank-forecast-exact}{20/20}
\setval{rank-flips-sel}{0}
\setval{rank-flips-alt}{10}
\setval{ledger-prereg}{8/8 and 0/6}
\setval{ledger-prereg-hit}{8 of 8}
\setval{ledger-prereg-fa}{0 of 6}
\setval{schema-fires}{0}
\setval{ledger-full-hits}{40/40}
\setval{ledger-full-fa}{0/30}
\setval{llm-explicit-exact}{31/31}
\setval{llm-under-diverge}{59/60}
\setval{llm-total-pipes}{91}
\setval{b3-crash-panels}{4}
\setval{minlen-l24-bank-range}{1.2--10.4\%}
\setval{minlen-l24-eia}{20.4\%}
\setval{minlen-l48-eia}{82.2\%}
\setval{seg-interp-max}{0.5\%}
\setval{abl-replication}{0.98--1.06}
\setval{prop-cases}{20}
\setval{prop-all-hold}{all}
\setval{prop4-cases}{50}
\setval{prop4-maxerr}{8e-16}
\setval{b1-eia-T1}{0.043}
\setval{b1-eia-bound2b}{0.149}
\setval{b1-bank-bound2b-max}{0.360}
\setval{mech-spread-none}{0.153}
\setval{mech-spread-tech}{0.092}
\setval{mech-spread-year}{0.119}
\setval{mech-reduction}{40\%}
\setval{mech-p-techyear}{0.049}
\setval{solar-q1}{44\%}
\setval{solar-q4}{63\%}
\setval{mech-b1-tech-share}{64\%}
\setval{tech-matched}{93\%}
\setval{d-inc-eia}{-0.57}
\setval{d-q12-eia}{-0.13}
\setval{d-q34-eia}{0.02}
\setval{e2-top1-share-bhc}{68\%}
\setval{e2-top1-share-bank}{30\%}
\setval{e2-b1-tv-bhc}{0.87}
\setval{e1-b1-tv-bhc}{0.41}
\setval{prereg-pass}{25}
\setval{prereg-fail}{23}
\setval{prereg-vacuous}{8}
\setval{prereg-total}{56}

\newcommand{\datatable}[1]{{\small\setlength{\tabcolsep}{4pt}\resizebox{\ifdim\width>\linewidth\linewidth\else\width\fi}{!}{\input{generated/#1}}}}
\begin{document}
\begin{frontmatter}
\title{How Calendar and History Filters Change Panel Forecast Comparisons}
\runtitle{Calendar and History Filters in Forecast Evaluation}
\author[1]{\inits{Y.}\fnms{Yunjie} \snm{Guo}}
\author[1]{\inits{H.}\fnms{Hanqing} \snm{Yang}\thanksref{c1}\ead{hanqing.yang189@gmail.com}}
\thankstext[type=corresp,id=c1]{Corresponding author.}
\address[1]{Independent Researcher, Suzhou, \cny{China}}
\begin{abstract}
Forecast benchmarks average losses over the series that pass preprocessing. A three-month grid and a calendar-quarter grid can therefore give different answers on the same records. We quantify this effect in four US banking panels and a panel of electricity generation. For a candidate that only removes whole series, reweighting the reference losses recovers the change in any weighted mean and in the ranking of separately fitted forecasters. A decomposition attributes the change to shifts between and within observed groups. The three-month grid removes 2,118 of 11,583 evaluated plants and widens the error advantage of the last-observation forecast over the historical mean by 17.3\% (paired entity-bootstrap interval, 13.3--21.1\%). Across banking panels the magnitude ranges from 0.8\% to 2.6\%. The absolute shift is comparable for plants and bank leverage, but the reference gap is much smaller for plants. Grouping plants by technology accounts for 64\% of the signed shift between groups. All four selection rules preserve the winning forecaster on every panel. Reporting retained counts alongside signed changes in loss reveals which benchmark conclusions depend on a preprocessing convention.
\end{abstract}
\begin{keyword}\kwd{forecast accuracy}\kwd{panel data}\kwd{sample selection}\kwd{standardization}\kwd{survivorship}\end{keyword}
\end{frontmatter}

\section{Introduction}
A bank first observed on 30 June reports again on 30 September and 31 December. Adding six months to its first date produces 30 December. A join on that date labels an observed quarter as missing. Repeat this operation and an otherwise complete series can fail a coverage screen. Banks that happen to start on 31 March follow a different path through the same code. Each surviving row can have a valid key, date and value, while the evaluation now answers a question about a different population.

This example is motivated by quarterly data preparation for the Financial Time Series Forecasting Repository \citep{ftsfr}. We study a controlled implementation in which the grid is the only change. The question is how much the composition of the evaluation sample contributes to the reported gap between two forecasts. A smaller sample can reduce coverage substantially while leaving that gap almost unchanged. Conversely, the gap can move materially while the ranking stays fixed. These possibilities matter when comparing results across benchmarks or deciding whether an analysis has to be rerun.

Existing work describes pitfalls in forecasting evaluation \citep{hewamalage}, automated checks for data quality \citep{schelter}, and the effects of selecting surviving funds on measured performance \citep{elton}. Our accounting uses the standard distinction between group composition and group-specific outcomes \citep{kitagawa}. Its application here requires two choices. First, we compare a candidate with a declared reference that already applies legitimate quality screens. Second, we hold the reference loss of each retained entity fixed to measure selection, then isolate any residual change requiring a new evaluation. The resulting calculation can be done for each pair of forecasters, even when selection leaves the identity of the winner unchanged.

\section{Panels, pipelines and outcomes}\label{sec:data}
The banking data are quarterly US Call Report series from 1976 through 2020, based on the consistent historical files of \citet{drechsler}. We examine assets divided by equity and cash divided by assets for individual banks and bank holding companies (BHCs). BHC ratios use sums of the component amounts of reporting banks within each holding company and date. The two variables and two aggregation levels provide sensitivity comparisons within one reporting system. The second domain is quarterly net electricity generation by US plant, from the EIA-923 series and its predecessor files for 2001--2024. We sum reported monthly generation within each plant and quarter and trim leading and trailing quarters without nonzero generation. Entry thus means first observed activity in this window. Of 15,131 plants in the constructed panel, \val{n-eia_plants} qualify for reference evaluation. All four banking panels and their denominators appear in Table~\ref{tab:grid}.

The reference grid runs between the first and last observation of each entity at calendar quarter ends. It permits at most 20\% missing quarters and drops the remaining missing observations. The three-month grid (B1) instead uses each first date plus an integer multiple of three months. Series that start in June or September miss subsequent March and December quarter ends. At the lengths used for evaluation, these series fail the coverage screen. For the other start months the grids coincide. We also consider a join restricted to the last snapshot, a rule requiring activity in the final year, a history-length filter, and alterations such as forward filling. Appendix~\ref{app:catalogue} gives the operations.

For entity \(i\) we hold out its last four available observations. We require at least 16 observations overall and a positive mean absolute first difference in the training series. Thus the four test observations need not span four consecutive calendar quarters when the retained history has gaps. Mean absolute scaled error (MASE) divides mean absolute test error by this training scale \citep{hyndman}. Our primary outcome \(a_i\) is MASE of the last-observation forecast minus MASE of the training mean. Negative values favor the last observation. The primary estimand, \(\theta\), gives equal weight to entities. We report signed shifts in MASE units and divide by \(\lvert\theta_r\rvert\) for relative shifts, where \(r\) denotes the reference. The study protocol calls a change material when its magnitude reaches 10\% of the reference gap. This cutoff is an interpretive convention and presupposes a nonzero reference gap.

Two secondary outcomes assess dependence on the estimand. We rank six entity-level forecasts by mean MASE: last observation, historical mean, seasonal naive (period four), drift, last-four mean, and simple exponential smoothing with its smoothing coefficient selected by training squared error. We also compute a pooled lag-one persistence statistic using the sum of within-entity lag products divided by the sum of squared centered lags, centering both series by each entity's full-series mean. For this statistic, the reference weights are squared-lag sums. The financial leverage ratios make those weights highly concentrated.

\section{Auditing a change in sample}\label{sec:method}
\subsection{Checking keys and retention}
The ledger compares candidate and reference processing against the same raw keys. It flags a candidate when a processed key has no raw counterpart, when its fraction of raw keys retained falls below half that of the reference, or when a retention spread exceeds the corresponding reference spread by more than 0.10. A spread is the maximum minus minimum retention rate across groups. Entity and key retention are checked by entry group and exit status. Key retention is also checked by age in years, capped at five. Entry group records both presence at the start of the panel and the calendar quarter of the first observation. Exit status records presence at the panel's final quarter.

An earlier ledger compared directly with raw retention. On a development subsample its flags included the benign reshuffle, because the reference screen itself changed retention. We fixed the differential version and logged predictions for fourteen variants before testing it on a 4,000-entity BHC-leverage subsample. The eight expected flags and six expected nonflags were all observed. The same pattern held on the five full panels. These counts measure agreement with specified test cases. They do not estimate sensitivity or specificity over naturally occurring pipeline errors. The prespecified schema check (columns, nonmissing values, unique keys, quarter-end dates) flagged none of the fourteen variants in the held-out subsample. Row-count and marginal-distribution checks can flag severe attrition, but do not identify which groups account for it.

\subsection{Attributing a change in the estimand}\label{sec:identity}
Let \(S_r\) and \(S_p\) be the evaluated reference and candidate sets. The reference has contributions \(a_i\) and positive weights \(w_i\). Let \(K=S_r\cap S_p\) have positive weight and let \(\theta_K\) be the reference weighted mean on \(K\). For a partition of \(S_r\), write \(W_{rc}\) and \(W_{kc}\) for normalized reference weight in cell \(c\) in the full and retained populations, respectively. Write \(\mu_{rc}\) and \(\mu_{kc}\) for the corresponding cell means, taking \(\mu_{kc}=\mu_{rc}\) in empty retained cells. Adding and subtracting \(\sum_c W_{kc}\mu_{rc}\) gives
\begin{equation}
\theta_p-\theta_r =
\underbrace{\sum_c(W_{kc}-W_{rc})\mu_{rc}}_{T_1}
+\underbrace{\sum_cW_{kc}(\mu_{kc}-\mu_{rc})}_{T_2}
+\underbrace{\theta_p-\theta_K}_{T_3}. \label{eq:decomp}
\end{equation}
The first term measures redistribution between cells. The second measures selection within cells. The residual includes changes to retained outcomes or weights and entities newly admitted by the candidate. The primary partition crosses entry group with exit status. Changing the partition redistributes \(T_1\) and \(T_2\) while preserving their sum. If no entity is retained, the ledger still reports the loss of coverage, but this decomposition is undefined.

For a candidate that only removes complete entities, \(K=S_p\). If retained series, weights and evaluation windows are unchanged and models are fit separately to each entity, then \(\theta_p=\theta_K\) and \(T_3=0\). A new set of eligible IDs and cached reference scores suffice to compute the change, including all pairwise mean-loss gaps. Refitting a global model on a changed training panel generally invalidates this argument. Removing dates within a retained history can also change its forecast or its scale.

Define total variation \(V=\tfrac12\sum_c|W_{kc}-W_{rc}|\). The sum of positive weight changes equals the absolute sum of negative changes, both equal to \(V\). Hence
\begin{equation}
 |T_1|\le V\,(\max_c\mu_{rc}-\min_c\mu_{rc}). \label{eq:bound}
\end{equation}
If weight moves between groups with equal mean losses the contribution vanishes. The bound concerns \(T_1\) alone. Substantial within-cell selection is possible even when the retained group shares match the reference. Appendix~\ref{app:bounds} gives the complementary bound for \(T_2\).

\section{Results}
\subsection{Calendar anchoring and the scale of the comparison}
\begin{table}[t]
\caption{Effect of the three-month grid on the last-observation minus historical-mean MASE gap. Relative shifts and their paired entity-bootstrap intervals are percentages of the absolute reference gap.}\label{tab:grid}
\centering\datatable{tab_grid}
\end{table}
The grid removes about two fifths of the evaluated holding companies and about one sixth of evaluated plants (Table~\ref{tab:grid}). The largest relative change is on plants: the gap moves from \(-0.246\) to \(-0.289\), increasing the advantage of the last-observation forecast. Bank leverage has a slightly larger absolute shift, \(-0.050\), from a reference gap of \(-2.979\). The small denominator on plants is therefore important to the cross-domain difference. The two BHC estimates move toward zero despite the greater attrition.

We added paired percentile intervals after the main experiments to quantify sensitivity to the mix of entities. Within each panel, we resampled reference entities with replacement 4,000 times, carrying each entity's reference loss and retention indicator together, and recomputed both means and the signed relative shift. The observed means and counts were checked against both original pipeline runs. The four banking intervals lie within the prespecified materiality band. The plant interval lies entirely below its lower boundary. At a 20\% threshold the plant point estimate would also fall below the cutoff. These are marginal intervals under exchangeable entity sampling. They do not capture common calendar shocks, geographic dependence, or links between banks in the same holding company. For the observed panels the shifts themselves are finite-population calculations.

\subsection{Why the plant sample changes}
For the grid rule, an entity's starting calendar quarter determines whether it is retained. With entry group included in the partition, the between-cell term accounts for the full shift on each panel. On plants the reference mean gap is \val{d-inc-eia} for incumbents. Among entrants it averages \val{d-q12-eia} for Q1--Q2 starts and \val{d-q34-eia} for Q3--Q4 starts. The grid excludes Q2 and Q3 entrants and keeps the incumbents. Both their loss profiles and their weights enter equation~\eqref{eq:decomp}.

We examined whether plant technology accounts for this pattern. The 2024 EIA-860 generator inventory supplies a technology for \val{tech-matched} of evaluated plants, using the largest generator by nameplate capacity and retaining an unmatched category. Technology labels outside the named categories are pooled into an additional other group. The solar share rises from \val{solar-q1} of Q1 entrants to \val{solar-q4} of Q4 entrants. Among entrants, adding technology indicators to a regression of the gap on entry quarter reduces the range of quarter coefficients from \val{mech-spread-none} to \val{mech-spread-tech}, a \val{mech-reduction} reduction against an advance prediction of at least 50\%. Entry-year indicators alone reduce it to \val{mech-spread-year}. With both technology and year, the joint quarter test gives \(p=\val{mech-p-techyear}\). This is a nominal descriptive test using an ordinary least-squares F statistic.

Grouping the full reference population by technology assigns \(-0.0271\) of the \(-0.0426\) shift to changing technology weights and \(-0.0155\) to selection within technologies. The between-technology share, \val{mech-b1-tech-share}, is an accounting result conditional on the partition. Technology recorded in 2024 is an imperfect proxy for a plant's technology at first observation. This exercise does not identify the causal effect of technology or season of entry.

\subsection{Other rules and outcomes}
\begin{table}[t]
\caption{Absolute relative change in the primary gap (\%). Bold values reach the 10\% materiality threshold. A superscript \(\ell\) indicates a ledger flag.}\label{tab:design}
\centering\datatable{tab_design}
\end{table}
Requiring membership in a final snapshot or activity in the final year changes the primary gap by \val{surv-bank-range} across banking panels and \val{surv-eia-range} on plants (Table~\ref{tab:design}). A minimum of 40 quarters produces large changes on plants and BHC cash but a 4.4\% change on bank leverage. For these three rules and the grid rule, restricting the reference scores recovers the rerun ranking in \val{rank-forecast-exact} comparisons. None changes the winning model. This result concerns six separately fitted forecasters at one horizon.

Forward filling a balanced panel also changes the trajectories supplied to forecasters. On plants its absolute shift exceeds 100\%: the magnitude of the original gap more than doubles in the same direction. The ledger detects the added keys but cannot measure this alteration from reference losses alone. In the secondary persistence statistic one entity accounts for \val{e2-top1-share-bhc} of total weight on BHC leverage and \val{e2-top1-share-bank} on bank leverage. The study therefore reports these values as a diagnostic of weight concentration and restricts substantive persistence comparisons to the cash panels. An audit should report the effective concentration of weights whenever it differs markedly from the fraction of entities retained.

\subsection{Minimum history as a sensitivity parameter}
Let \(n\) be the number of observations in the reference series. Raising the evaluation minimum from \(L_0=16\) to \(L>L_0\) removes some short histories. With \(\pi_L\) the retained reference weight and \(\mu_{\ge L},\mu_{<L}\) the two group means, conditioning on the filter gives
\begin{equation}
\theta(L)-\theta_r=(1-\pi_L)(\mu_{\ge L}-\mu_{<L}). \label{eq:length}
\end{equation}
If the retained set is unchanged the shift is zero. For positive mass on each side, both factors can be estimated directly from the reference scores. With monotone weighted conditional means by length, the retained tail mean is monotone as well. Monotonicity is an empirical question in other panels.

Figure~\ref{fig:length} uses a second, task-specific reference that includes a synthetic point-in-time eligibility attribute independent of the outcome. It varies only the length filter, retaining the same eligible observations and evaluator. It therefore has a different denominator from Table~\ref{tab:design}. Requiring 24 quarters changes the gap by \val{minlen-l24-bank-range} on banking panels and by \val{minlen-l24-eia} on plants. At 48 quarters the change reaches \val{minlen-l48-eia} on plants. Across 50 panel--threshold comparisons, the analytic and rerun shifts agree to machine precision. The banking sweep was exploratory. Its monotone direction was predicted in advance for plants.
\begin{figure}[t]
\centering
\begin{tikzpicture}
\begin{axis}[width=.94\linewidth,height=6cm,xmin=15,xmax=49,ymin=0,ymax=90,
xlabel={Required observations per entity},ylabel={Absolute relative shift (\%)},
xtick={16,20,24,28,32,40,48},ytick={0,20,40,60,80},ymajorgrids,grid style={gray!25},
legend style={at={(0.015,.98)},anchor=north west,font=\footnotesize,draw=none,fill=white},legend columns=2]
\addplot[thick,blue!65!black,mark=square*,mark size=1.3pt] table[x=L,y=bank_leverage]{generated/fig_minlen.dat};\addlegendentry{Bank leverage}
\addplot[thick,teal!75!black,mark=triangle*,mark size=1.6pt] table[x=L,y=bank_cash_liquidity]{generated/fig_minlen.dat};\addlegendentry{Bank cash}
\addplot[thick,blue!65!black,densely dashed,mark=square,mark size=1.5pt] table[x=L,y=bhc_leverage]{generated/fig_minlen.dat};\addlegendentry{BHC leverage}
\addplot[thick,teal!75!black,densely dashed,mark=triangle,mark size=1.6pt] table[x=L,y=bhc_cash_liquidity]{generated/fig_minlen.dat};\addlegendentry{BHC cash}
\addplot[very thick,red!60!black,mark=*,mark size=1.6pt] table[x=L,y=eia_plants]{generated/fig_minlen.dat};\addlegendentry{Power plants}
\end{axis}
\end{tikzpicture}
\caption{Effect of increasing minimum history, relative to the gap at 16 observations in the task-specific reference. Markers show evaluated thresholds and lines connect adjacent measurements.}\label{fig:length}
\end{figure}

\section{Discussion}
For a forecaster deciding whether to rerun an evaluation, the first question is whether the candidate changes retained training or test values. If it only removes whole series and each model fits locally, publishing per-entity losses and the membership indicator is enough to recover the new mean and ranking. Reporting relative change beside the signed change and both denominators makes the result interpretable when the original gap is small. For a change in values or model training, direct evaluation is needed to measure the residual.

The reference specifies the target population; accuracy on a population absent from the raw files remains unidentified. A ledger flag describes a difference even when the candidate implements a legitimate restriction, and uniform losses in unmeasured groups can escape its spread check. The four financial panels share reporting institutions, and the electricity data provide only one cross-domain replication. The design measures sensitivity to specific conventions on these five panels. It does not estimate the prevalence of such conventions in forecasting practice.

The repository preserves the dated protocol, code hashes, deviations and unsuccessful predictions. These are project-maintained records, rather than an independently time-stamped registry. The banking and electricity panels contradicted opposite predictions about the grid rule: the development work anticipated a larger effect on cash ratios, while the electricity study predicted a sub-threshold shift. In both settings the observable membership changes and reference losses suffice to explain the realized direction and magnitude.

\section*{Data and code availability}
The replication materials include scored entity-level data for the calendar comparison, the paired resampling script, original experiment outputs, panel construction details and the dated protocol. The EIA source manifest lists annual files and hashes. The README documents reconstruction from the NYU Call Reports and EIA public releases. The version accompanying this manuscript is maintained at \url{https://github.com/fightforbelief/p18-convention-errors}.

\begin{appendix}
\section{Bounds for the selection terms}\label{app:bounds}
Using the reference distribution within each cell, let \(k_i\) denote retention, \(\pi_c=E_c[k]\) its weighted mean and \(\sigma_c\) the weighted standard deviation of \(a_i\). For a nonempty retained cell,
\[
\mu_{kc}-\mu_{rc}=\operatorname{Cov}_c(a,k)/\pi_c.
\]
Cauchy--Schwarz and \(\operatorname{Var}_c(k)=\pi_c(1-\pi_c)\) imply
\[
 |T_2|\le \sum_{c:\pi_c>0}W_{kc}\sigma_c\sqrt{(1-\pi_c)/\pi_c}.
\]
If each cell is fully retained or removed, then \(T_2=0\). Write \(\sigma_B^2=\sum_c W_{rc}(\mu_{rc}-\theta_r)^2\) for the between-cell variance. Subtract \(\theta_r\) from each mean in \(T_1\), since the changes in weights sum to zero. Cauchy--Schwarz then gives
\[
 |T_1|\le \sigma_B\left\{\sum_{c:W_{rc}>0}W_{rc}(W_{kc}/W_{rc}-1)^2\right\}^{1/2}.
\]
The bounds need no assumptions about why an entity was retained. They apply to reference-weight cells and can be loose when changes in cell contributions offset each other.

\section{Pipeline catalogue}\label{app:catalogue}
The B variants were used in the design of the audit. The x variants were authored after the first ledger was frozen and draw on known pitfalls. The differential ledger was developed subsequently. These are controlled implementations. In x8, ``benign'' describes its effect on the scale-invariant forecast gap in this design.
\begin{center}\datatable{tab_catalogue}\end{center}
\end{appendix}
\bibliographystyle{jds2}
\bibliography{references}
\end{document}
